%%=============================================================================
%% Samenvatting
%%=============================================================================

% TODO: De "abstract" of samenvatting is een kernachtige (~ 1 blz. voor een
% thesis) synthese van het document.
%
% Een goede abstract biedt een kernachtig antwoord op volgende vragen:
%
% 1. Waarover gaat de bachelorproef?
% 2. Waarom heb je er over geschreven?
% 3. Hoe heb je het onderzoek uitgevoerd?
% 4. Wat waren de resultaten? Wat blijkt uit je onderzoek?
% 5. Wat betekenen je resultaten? Wat is de relevantie voor het werkveld?
%
% Daarom bestaat een abstract uit volgende componenten:
%
% - inleiding + kaderen thema
% - probleemstelling
% - (centrale) onderzoeksvraag
% - onderzoeksdoelstelling
% - methodologie
% - resultaten (beperk tot de belangrijkste, relevant voor de onderzoeksvraag)
% - conclusies, aanbevelingen, beperkingen
%
% LET OP! Een samenvatting is GEEN voorwoord!

%%---------- Nederlandse samenvatting -----------------------------------------
%
% TODO: Als je je bachelorproef in het Engels schrijft, moet je eerst een
% Nederlandse samenvatting invoegen. Haal daarvoor onderstaande code uit
% commentaar.
% Wie zijn bachelorproef in het Nederlands schrijft, kan dit negeren, de inhoud
% wordt niet in het document ingevoegd.

\IfLanguageName{english}{%
\selectlanguage{dutch}
\chapter*{Samenvatting}
\selectlanguage{english}
}{}

%%---------- Samenvatting -----------------------------------------------------
% De samenvatting in de hoofdtaal van het document

\chapter*{\IfLanguageName{dutch}{Samenvatting}{Abstract}}

Onvoorziene defecten van bouwmachines kunnen resulteren in stilstand, vertragingen en extra onderhoudskosten. Voorspellend onderhoud creëert kansen om afwijkend machinegedrag op tijd te waarnemen en onderhoud effectiever te organiseren. Voor kleine bouwbedrijven kan dit echter een uitdaging zijn, omdat er vaak slechts een beperkte hoeveelheid historische sensordata en een beperkt aantal geregistreerde defecten beschikbaar is. Deze bachelorproef bestudeert hoe een machine learning-methodiek, die steunt op sensordata, kan worden ingezet om storingen in bouwmachines in een kleine bouwbedrijf te voorspellen en het onderhoud efficiënter te plannen.

Voor het onderzoek zijn real-world telematicadata van een tractor die wordt ingezet voor grondverzet en transport gebruikt. Eerst werden de data opgeslagen, samengevoegd en voorbereid als een tijdsreeksdataset. Daarna werden belangrijke sensoren gekozen die informatie verstrekken over onder andere temperatuur, druk, motorbelasting, motortoerental en brandstofverbruik. Omdat de private tractordataset geen duidelijke defect- of onderhoudslabels had, werd er extra gebruikgemaakt van de openbare MetroPT-3-dataset. Deze dataset omvat industriële tijdsreeksdata met vastgelegde defectmomenten, waardoor het mogelijk werd  om de ontwikkelde methodiek objectiever te beoordelen.

In de methodiek zijn vijf modellen bestudeerd: Random Forest-model, Isolation Forest-model, Autoencoder, LSTM Autoencoder en One-Class SVM-model. De modellen werden aanvankelijk gebruikt op de MetroPT-3-dataset en daarna, waar mogelijk, op de private tractordataset. De uitkomsten gaven aan dat de prestaties aanzienlijk afhankelijk zijn van de gekozen methode en evaluatiemaatstaf. De Autoencoder leverde de beste resultaten op voor de defectklasse tijdens de aangepaste MetroPT-3-testperiode met een recall van 0.44 en een F1-score van 0.24. Random Forest-model en Isolation Forest-model ontdekten slechts een klein aantal van de geregistreerde defecten, terwijl One-Class SVM-model meer defecten kon opsporen, maar tegelijkertijd veel gebruikelijke observaties als een anomalie beschouwde.

Het onderzoek laat zien dat machine learning kan worden gebruikt om afwijkend gedrag in industriële sensordata te identificeren, wat een fundament kan leggen voor voorspellend onderhoud van bouwmachines. De implementatie van de private tractordataset heeft aangetoond dat de methodiek ook toepasbaar is voor real-world sensorgegevens, maar dat de ontdekte afwijkingen door het ontbreken van defectlabels niet als bevestigde storingen kunnen worden gezien. De ontwikkelde methode fungeert daarom voornamelijk als een bewijs. Voor een betrouwbare praktijktoepassing zijn uitgebreide historische datasets,  verschillende onderhouds- en storingsmomenten vereist. Dit maakt het mogelijk om toekomstige modellen effectiever te valideren en verder te optimaliseren voor gebruik in kleine bouwbedrijven.

%Onvoorziene defecten van bouwmachines kunnen resulteren in stilstand, vertragingen en extra onderhoudskosten. Voor kleine bouwbedrijven biedt voorspellend onderhoud een boeiende kans om machineproblemen vroegtijdig te herkennen en het onderhoud effectiever te plannen. In de praktijk hebben dergelijke bedrijven vaak een beperkte hoeveelheid historische sensordata en zijn er weinig of geen geregistreerde defecten. Dit maakt het niet vanzelfsprekend om traditionele machine learning-modellen voor het voorspellen van storingen direct in te zetten. Deze bachelorproef bestudeert hoe een machine learning-methodiek, die steunt op sensordata, kan worden ingezet om storingen in bouwmachines in een kleine bouwbedrijf te voorspellen en het onderhoud efficiënter te plannen.

%Het onderzoek vond plaats in samenwerking met een kleine bouwbedrijf en gebruikte real-world telematicadata van een tractor die wordt gebruikt voor grondverzet en transport. Eerst werden de beschikbare data opgeslagen en voorbereid voor een verdere analyse. Onder andere zijn er diverse databestanden samengevoegd, zijn tijdstempels gestandaardiseerd, zijn onnodige data verwijderd en zijn sensordata omgezet naar de juiste numerieke waarden. Daarna werden belangrijke parameters gekozen die informatie verstrekken over onder andere motorbelasting, temperatuur, druk, motortoerental en brandstofverbruik.

%Een cruciale beperking van de private tractordataset was het gebrek aan heldere historische defect- en onderhoudslabels. Hierdoor was het niet mogelijk om direct vast te stellen of een gedetecteerde afwijking overeenkwam met een echte machinefout. Om de ontwikkelde methodiek op een dataset met bekende tekortkomingen te kunnen beoordelen, werd er extra gebruikgemaakt van de openbare MetroPT-3-dataset. Deze dataset omvat tijdsreeksgegevens van een industrieel compressorsysteem en registreert defectmomenten. hierdoor was het mogelijk om de effectiviteit van de diverse methoden te evalueren op basis van bekende afwijkingen.





